\begin{table}[H]
    \centering
    \caption{Predictors of ``I don't know'' responses.}
    \begin{tabular}{ l c c }
\hline \hline
 & estimate & se \\
\hline
\hspace{0.05cm} Age: 35-54 (vs. 18-34) & 0.0397** & (0.0199) \\
\hspace{0.05cm} Age: 55+ (vs. 18-34) & 0.0474** & (0.0206) \\
\hspace{0.05cm} Gender: Male (vs. Female) & -0.0691*** & (0.0134) \\
\hspace{0.05cm} Region: Northeast (vs. Northwest) & -0.0060 & (0.0198) \\
\hspace{0.05cm} Region: Center (vs. Northwest) & -0.0039 & (0.0197) \\
\hspace{0.05cm} Region: South (vs. Northwest) & 0.0071 & (0.0185) \\
\hspace{0.05cm} Region: Islands (vs. Northwest) & -0.0001 & (0.0244) \\
\hspace{0.05cm} Education: Upper secondary (vs. Compulsory/lower) & -0.0414 & (0.0269) \\
\hspace{0.05cm} Education: Tertiary (vs. Compulsory/lower) & -0.0971*** & (0.0282) \\
\hspace{0.05cm} Task: 1/FEW (vs. Task 2/ANY) & 0.0493*** & (0.0075) \\
\hspace{0.05cm} Task: 3/ANY2 (vs. Task 2/ANY) & 0.0049 & (0.0070) \\
\hline
\hspace{0.05cm} Observations & \multicolumn{2}{c}{3,245} \\
\hspace{0.05cm} R-squared & \multicolumn{2}{c}{0.015} \\
\hline \hline
\end{tabular}
\begin{tablenotes}
\small
\item \textit{Notes:} Sociodemographic characteristics and task/scenario features. Pooled linear probability model across the 3 DCE tasks, cluster-robust respondent-level SE in parentheses. *** p$<$0.01, ** p$<$0.05, * p$<$0.1. Baselines: Age vs. 18-34; Gender vs. Female ("Other", N=4, excluded); Region vs. Northwest; Education vs. Compulsory/lower secondary; Task vs. Task 2 (ANY). Political orientation not yet available (placeholder pending wave 1 merge).
\end{tablenotes}
    \label{tab:table12}
\end{table}
